\chuong{KIỂM CHỨNG PROOF OF CONCEPT}
\begingroup\setlength{\emergencystretch}{2em}
\subsection{Kịch bản kiểm chứng}

Mỗi yêu cầu ở Bảng~\ref{tab:yeucau} được gắn với ít nhất một kịch bản, có tiêu chí đạt viết sẵn từ trước và đầu ra kiểm tra được (Bảng~\ref{tab:scenarios}). Các kịch bản chạy ở hai mức. Ở mức tĩnh, kịch bản K9 đối chiếu ma trận kỳ vọng với danh mục công cụ, chính sách phân quyền và cấu hình realm mà không cần đến cụm. Ở mức đầu cuối, hai bộ kiểm thử, một cho danh mục công cụ và một cho phân quyền theo agent, dùng MCP Python SDK 1.30.0 gọi thật vào \texttt{/mcp} và dừng ngay ở phép so sánh sai đầu tiên. Kết quả ``PASS'' vì thế chỉ xuất hiện khi mọi phép so đều đúng. Riêng thăm dò NetworkPolicy và đọc access log được làm thủ công, đầu ra lưu lại làm bằng chứng.

\begin{table}[H]\centering\caption{Kịch bản kiểm chứng theo yêu cầu}\label{tab:scenarios}
{\singlespacing\renewcommand{\arraystretch}{1.15}\small
\begin{tabularx}{\textwidth}{|c|>{\raggedright\arraybackslash}p{1.5cm}|>{\raggedright\arraybackslash}X|>{\raggedright\arraybackslash}p{4.2cm}|}\hline
\textbf{Mã} & \textbf{Yêu cầu} & \textbf{Kịch bản và tiêu chí đạt} & \textbf{Bằng chứng} \\ \hline
K1 & F1, F2 & Liệt kê và gọi cả 21 công cụ qua \texttt{/mcp}; danh sách tên đã sắp xếp phải trùng khớp và không lệnh nào lỗi & \path{gateway-smoke.txt} (08/10) \\ \hline
K2 & F3, S2 & Mười agent xin token; \texttt{tools/list} của mỗi agent phải trùng ma trận kỳ vọng & \path{agent-rbac-smoke.txt} (29/10) \\ \hline
K3 & S1 & Mỗi agent gọi một công cụ ngoài quyền và phải bị từ chối; ẩn tạm một công cụ rồi gọi thẳng tên của nó & \path{agent-rbac-smoke.txt}, \path{tool-denial*.txt} \\ \hline
K4 & S2, S3 & Token mang đúng \texttt{aud}, \texttt{azp}, \texttt{roles}; backend báo lại đúng header danh tính & \path{agent-rbac-smoke.txt} \\ \hline
K5 & S1, S3 & Gọi bằng curl khi thiếu token, khi gửi header danh tính giả, và kiểm tra token gốc không tới backend & Chưa lưu đầu ra \\ \hline
K6 & S4 & Thăm dò backend radar từ namespace \texttt{default} và từ namespace của gateway & \path{agent-rbac-smoke.txt} \\ \hline
K7 & S5, O2 & Mọi lệnh cho phép và bị chặn phải có bản ghi kèm danh tính agent trong access log và trên dashboard & Access log, ảnh dashboard \\ \hline
K8 & O1 & Năm lần khởi động ấm liên tiếp, mỗi lần kèm kiểm thử phân quyền & \path{benchmarks/warm.md} \\ \hline
K9 & F2, S2 & Kiểm thử tĩnh đối chiếu ma trận, danh mục, chính sách và realm & \path{static_contract_test.mjs} \\ \hline
\end{tabularx}}\end{table}

Môi trường thử là cụm kind một node chạy Kubernetes v1.37.0 với Calico v3.33.0 trên WSL2, cùng agentgateway v1.5.0 và Keycloak 26.2.5. Bằng chứng được ghi trong ba đợt. Đợt ngày 08/10/2026 dành cho danh mục và thí nghiệm ẩn công cụ, đợt 29/10/2026 cho phân quyền theo agent, cô lập mạng và kiểm toán, còn đợt 05/11/2026 cho dashboard và thời gian khởi động. Các mục dưới đây trình bày kết quả lần lượt theo chức năng, phân quyền, an toàn rồi vận hành; Bảng~\ref{tab:ch5-goals} ở cuối chương gom lại theo từng yêu cầu.

\subsection{Kết quả chức năng}

Kịch bản K1 đạt trọn vẹn. Gateway trả về 21 công cụ, và cả 21 lệnh gọi đều thành công. Bộ kiểm thử so toàn bộ danh sách tên đã sắp xếp thay vì chỉ đếm, nên chỉ cần một công cụ thiếu, thừa hay sai tiền tố là lần chạy thất bại. Số công cụ theo target khớp Bảng~\ref{tab:tools}, với radar 4, spark-history 3, openmetadata 4, dbt 5 và airflow 5. Access log của cùng lượt chạy cho thấy gateway tách tên đầy đủ trở lại thành target và tên gốc, chẳng hạn \texttt{mcp.target=spark-history} đi cùng \texttt{gen\_ai.tool.name=get\_stage}. Đây cũng chính là hai trường mà quy tắc phân quyền đọc.

Có điều kiện đi kèm kết quả này. Lượt chạy ngày 08/10/2026 diễn ra trước khi bật xác thực JWT, nên access log ghi \texttt{user="anonymous"}. Sau khi bật JWT, danh mục đủ 21 công cụ được xác nhận lại qua \texttt{platform-admin-agent} ở kịch bản K2, nhưng agent này chỉ gọi một công cụ. F1 và F2 vì vậy có bằng chứng đầy đủ ở bước khám phá dưới JWT, trong khi bằng chứng gọi đủ 21 công cụ thuộc về cấu hình chưa bật JWT.

\subsection{Kết quả phân quyền}

Trong lượt chạy ngày 29/10/2026, cả mười agent đều đạt kịch bản K2. Bảng~\ref{tab:rbac} đặt số kỳ vọng, tính từ ma trận agent, vai trò và công cụ, cạnh số thực tế. Bộ kiểm thử so khớp từng tên chứ không so số lượng, nên việc hai cột trùng nhau có nghĩa là cả 210 ô của ma trận (10 agent nhân 21 công cụ) đều đúng ở bước khám phá.

\begin{table}[H]\centering\caption{Kết quả kiểm thử phân quyền theo agent}\label{tab:rbac}
{\singlespacing\renewcommand{\arraystretch}{1.15}\small
\begin{tabularx}{\textwidth}{|p{3.6cm}|X|>{\centering\arraybackslash}p{2.4cm}|>{\centering\arraybackslash}p{2.4cm}|}\hline
\textbf{Agent} & \textbf{Vai trò} & \textbf{Được phép} (kỳ vọng/ thực tế) & \textbf{Bị chặn} (kỳ vọng/ thực tế) \\ \hline
cluster-observer & kube\_viewer & 3/3 & 18/18 \\ \hline
incident-response & kube\_operator, spark\_viewer & 7/7 & 14/14 \\ \hline
spark-diagnostic & spark\_viewer & 3/3 & 18/18 \\ \hline
metadata-discovery & metadata\_viewer & 3/3 & 18/18 \\ \hline
metadata-steward & metadata\_editor & 4/4 & 17/17 \\ \hline
dbt-quality & dbt\_operator & 5/5 & 16/16 \\ \hline
airflow-monitor & airflow\_viewer & 3/3 & 18/18 \\ \hline
airflow-operator & airflow\_operator & 5/5 & 16/16 \\ \hline
platform-admin & platform\_admin & 21/21 & 0/0 \\ \hline
no-access & (không có) & 0/0 & 21/21 \\ \hline
\textbf{Tổng (210 ô)} & & \textbf{54/54} & \textbf{156/156} \\ \hline
\end{tabularx}}\end{table}

Hai dòng trong bảng kiểm tra trực tiếp hai tính chất của mô hình quyền. \texttt{incident-response-agent}, agent duy nhất mang hai vai trò, nhận đúng 7 công cụ, bằng 4 công cụ radar cộng 3 công cụ Spark; đó là quyền hợp mà S2 đặt ra. \texttt{no-access-agent} thấy danh mục trống dù không có quy tắc từ chối nào, bởi mọi quy tắc trong chính sách đều là quy tắc cho phép và đòi một vai trò cụ thể. Mặc định từ chối của S1 thể hiện đúng ở dòng này.

Bước gọi (K3) có độ phủ hẹp hơn. Với mỗi agent, bộ kiểm thử gọi công cụ đứng đầu theo thứ tự chữ cái trong danh sách được phép, rồi gọi công cụ đứng đầu ngoài quyền. Mười agent do đó sinh ra 18 lệnh, gồm 9 lệnh cho phép (trừ \texttt{no-access}) và 9 lệnh bị chặn (trừ \texttt{platform-admin}). Bảy trong số các lệnh bị chặn cùng nhắm vào \texttt{airflow\_get\_dag\_run}. Lệnh \texttt{airflow\_pause\_dag} của \texttt{airflow-monitor} là trường hợp duy nhất thử tách quyền đọc khỏi quyền ghi trên cùng một hệ thống bằng lệnh gọi thật.

Cơ chế từ chối có hai lớp, khớp với Hình~\ref{fig:decision}. Ở bước khám phá, công cụ ngoài quyền không xuất hiện trong \texttt{tools/list}. Còn khi một lệnh \texttt{tools/call} gửi thẳng tên của công cụ bị ẩn, gateway trả HTTP 400 kèm \texttt{Unknown tool}. Thí nghiệm ngày 08/10/2026 tách riêng được lớp thứ hai này. Nhóm tạm rút quyền với \texttt{airflow\_pause\_dag}, danh mục còn lại 20 công cụ, và lệnh gọi trực tiếp nhận \texttt{400 Bad Request}; access log ghi \texttt{reason=MCP} cùng \texttt{duration=0ms}, chứng tỏ gateway tự trả lỗi mà không chuyển yêu cầu xuống backend.

\subsection{Kiểm chứng an toàn}

\subsubsection{Token và danh tính}

Chiều chấp nhận của xác thực đã có bằng chứng (K4). Token của cả 10 agent mang đúng \texttt{aud}, \texttt{azp} và \texttt{roles} theo ma trận. Ở 9/9 lệnh được phép, mock báo lại rằng nó nhận \texttt{x-agent-id} bằng client ID của agent và \texttt{x-forwarded-user} bằng tên tài khoản dịch vụ tương ứng, tức hai giá trị lấy từ JWT đã kiểm tra chứ không phải do client tự điền.

Chiều từ chối thì chưa có bằng chứng được lưu (K5). Ba trường hợp, gồm mã 401 khi thiếu token, việc ghi đè header danh tính giả và việc token gốc không tới được backend, đã được chạy bằng curl nhưng đầu ra không được giữ lại, nên báo cáo không tính chúng là kết quả đã kiểm chứng. Token sai chữ ký, sai audience hoặc hết hạn vẫn chưa có kịch bản. S3 thì có thể khép lại với chi phí rất nhỏ. Mock vốn đã báo về trường \texttt{authorization\_present}, chỉ là bộ kiểm thử phân quyền chưa khẳng định trường này bằng \texttt{false}; thêm một phép so là có bằng chứng tự động cho cả 9 lệnh được phép.

\subsubsection{Cô lập mạng}

Kịch bản K6 gồm hai lần thăm dò tới backend radar, và hai lần này phải đi cùng nhau mới có giá trị. Kết nối từ namespace \texttt{default} hết thời gian chờ. Kết nối từ namespace của gateway lại nhận HTTP 406 từ chính MCP endpoint, mã mà endpoint Streamable HTTP trả về khi yêu cầu không khai báo nhận được cả JSON lẫn SSE, đúng như một lần thăm dò thô. Mã 406 do đó chứng tỏ backend vẫn chạy và đã trả lời. Nếu chỉ có lần thăm dò đầu, việc hết thời gian chờ hoàn toàn có thể do backend đang tắt chứ không phải do NetworkPolicy. Ghép hai kết quả lại mới thấy backend chỉ bị chặn theo nguồn (mục~\ref{sec:netpol}). Backend Docker chạy ngoài cụm chưa được thăm dò.

\endgroup
\subsection{Kiểm chứng tự động hóa và observability}

Ngày 05/11/2026, nhóm chạy một bộ đo lặp năm lần liên tiếp trên máy chủ thử nghiệm, khi cụm và mọi image đã có sẵn, tức tình huống khởi động hằng ngày. Mỗi lần gồm ba bước nối tiếp, lần lượt là áp toàn bộ các lớp cấu hình và cài bộ giám sát, mở kết nối tới gateway, Keycloak và Grafana rồi chờ chúng phản hồi, và cuối cùng chạy trọn bộ kiểm thử phân quyền cho mười agent. Sau bước cuối, số lệnh cho phép và bị chặn được đếm trực tiếp trong access log của gateway. Bảng~\ref{tab:bench} tổng hợp kết quả.

\begin{table}[H]\centering\caption{Kết quả đo khởi động ấm và kiểm thử phân quyền qua năm lần chạy}\label{tab:bench}
{\singlespacing\renewcommand{\arraystretch}{1.15}\small
\begin{tabularx}{\textwidth}{|X|>{\centering\arraybackslash}p{1.6cm}|>{\centering\arraybackslash}p{1.6cm}|>{\centering\arraybackslash}p{1.6cm}|>{\centering\arraybackslash}p{1.9cm}|}\hline
\textbf{Chỉ số (n = 5)} & \textbf{Nhỏ nhất} & \textbf{Trung bình} & \textbf{Lớn nhất} & \textbf{Độ lệch chuẩn} \\ \hline
Áp cấu hình và cài giám sát (giây) & 4,114 & 4,148 & 4,199 & 0,034 \\ \hline
Đến khi các dịch vụ phản hồi (giây) & 5,142 & 5,177 & 5,230 & 0,036 \\ \hline
Kiểm thử phân quyền 10 agent (giây) & 1,864 & 1,896 & 1,931 & 0,026 \\ \hline
Lệnh cho phép / bị chặn trong access log & 9 / 9 & 9 / 9 & 9 / 9 & 0 \\ \hline
Lần kiểm thử phân quyền đạt & \multicolumn{4}{c|}{5/5} \\ \hline
\end{tabularx}}\end{table}

Thời gian giữa các lần chênh nhau chưa tới 0,09 giây, độ lệch chuẩn dưới 1\% giá trị trung bình, nên kết quả ổn định và lặp lại được. Khoảng chênh khoảng 1,03 giây giữa hai dòng thời gian đầu là phần dành cho việc mở kết nối và chờ dịch vụ phản hồi, gần như không đổi qua các lần. Tính từ lúc bắt đầu, hệ thống sẵn sàng phục vụ agent sau khoảng 5,2 giây, và một lượt kiểm chứng toàn bộ ma trận quyền mất thêm chưa tới 2 giây. Vì mọi tài nguyên đã tồn tại, các lệnh \texttt{kubectl apply} chỉ xác nhận trạng thái không đổi, realm không bị tạo lại và secret của agent được giữ nguyên. Năm lần chạy liên tiếp đều cho kết quả phân quyền đạt chính là bằng chứng cho tính idempotent, vì chạy lại script trên hệ thống đang hoạt động không đổi danh tính hay khóa ký, nên token đang dùng vẫn hợp lệ.

Ở cả năm lần, access log ghi đúng 9 lệnh cho phép và 9 lệnh bị chặn, khớp dự đoán từ ma trận, gồm chín agent có công cụ được phép (trừ \texttt{no-access}) và chín agent có công cụ bị cấm (trừ \texttt{platform-admin}). Hình~\ref{fig:allowdeny} cho thấy dashboard Grafana hiển thị cùng số liệu cho một lượt chạy. Ở các ô chi tiết (Hình~\ref{fig:dashboard}), lệnh bị chặn của \texttt{no-access-agent} xuất hiện với \texttt{agent\_roles=[]}, \texttt{http.status=400} và \texttt{Unknown tool}, còn lệnh của \texttt{platform-admin-agent} trên cùng công cụ ghi \texttt{http.status=200}. Chỉ từ dashboard, người vận hành trả lời được agent nào gọi, gọi công cụ nào và kết quả ra sao, trong khi token không xuất hiện trong log.

\begin{figure}[H]\centering
\includegraphics[width=0.4\textwidth]{grafana-allow.png}\hspace{1em}
\includegraphics[width=0.4\textwidth]{grafana-deny.png}
\caption{Số lệnh cho phép và bị chặn trên dashboard sau một lượt kiểm thử}\label{fig:allowdeny}\end{figure}

\begingroup\setlength{\emergencystretch}{2em}
\subsection{Đánh giá}
\subsubsection{Đối chiếu mục tiêu}

Bảng~\ref{tab:ch5-goals} tổng hợp số đo của từng yêu cầu. Một yêu cầu được xếp ``Đạt'' khi mọi phép đo đã chạy cho kết quả bằng kỳ vọng và phần cốt lõi có đầu ra đã lưu. Mức ``Một phần'' dành cho trường hợp các phép đo đã chạy đều đạt nhưng vẫn còn tình huống chưa có đầu ra hoặc nằm ngoài phạm vi thử. Tiêu chí này cố ý đặt chặt, đến mức một kết quả đã quan sát được mà không lưu đầu ra cũng không được tính.

\begin{table}[H]\centering\caption{Tổng hợp kết quả đo theo yêu cầu}\label{tab:ch5-goals}
{\singlespacing\renewcommand{\arraystretch}{1.15}\small
\begin{tabularx}{\textwidth}{|c|X|>{\centering\arraybackslash}p{1.7cm}|>{\centering\arraybackslash}p{1.7cm}|>{\centering\arraybackslash}p{1.9cm}|}\hline
\textbf{Mã} & \textbf{Phép đo} & \textbf{Kỳ vọng} & \textbf{Thực tế} & \textbf{Mức đạt} \\ \hline
F1 & Lệnh gọi thành công qua \texttt{/mcp} & 21 & 21 & Đạt \\ \hline
F2 & Tên công cụ đúng tiền tố target & 21 & 21 & Đạt \\ \hline
F3 & Ô ma trận khớp ở \texttt{tools/list} & 210 & 210 & Đạt \\ \hline
S1 & Lệnh ngoài quyền bị chặn & 9 & 9 & Một phần \\ \hline
S2 & Agent có claim \texttt{roles} đúng & 10 & 10 & Đạt \\ \hline
S3 & Lệnh mà backend nhận token gốc & 0 & chưa lưu & Một phần \\ \hline
S4 & Lần thăm dò mạng cho kết quả đúng & 2 & 2 & Một phần \\ \hline
S5 & Lệnh gọi có bản ghi kiểm toán & 18 & 18 & Đạt \\ \hline
O1 & Lần khởi động ấm đạt kiểm thử (TB 4,15 giây) & 5 & 5 & Một phần \\ \hline
O2 & Lần đếm allow / deny khớp 9 / 9 & 5 & 5 & Đạt \\ \hline
\end{tabularx}}\end{table}

Sáu yêu cầu đạt, bốn yêu cầu đạt một phần. Ở mọi phép đo đã chạy, giá trị thực tế đều bằng kỳ vọng, và không phép đo nào cho kết quả sai. Bốn yêu cầu đạt một phần cùng chung nguyên nhân là thiếu đầu ra hoặc phạm vi thử hẹp. S1 chưa có đầu ra cho mã 401 và token sai, S3 chưa có đầu ra cho việc xóa token gốc, S4 chưa thăm dò backend ngoài cụm, O1 chưa đo lần dựng từ cụm trống. Mục tiêu chung ở Chương 2 như vậy đạt được ở phần lõi, tức danh mục hợp nhất, phân quyền theo công cụ và cô lập mạng trong cụm. Phần chưa đạt lại rơi vào chiều từ chối của xác thực, cũng là chiều mà kẻ tấn công sẽ thử trước tiên.

\subsubsection{Điểm mạnh}

Điểm mạnh dễ thấy nhất là ma trận kỳ vọng được viết thành dữ liệu và so khớp chính xác. Cùng một ma trận vừa là kết quả mong đợi của kiểm thử đầu cuối, vừa là đầu vào của kiểm thử tĩnh (K9), và lần chạy lại kiểm thử tĩnh ngày 05/11/2026 vẫn đạt với 21 công cụ và 10 agent mà không cần cụm. Bất kỳ thay đổi chính sách nào không kèm cập nhật ma trận sẽ bị bắt ngay trên máy cá nhân, trước khi kịp áp vào cụm.

Điểm mạnh thứ hai nằm ở chỗ việc từ chối diễn ra sớm và ngay tại gateway. 156 trên 210 ô bị ẩn trước khi mô hình kịp chọn công cụ, còn lệnh cố tình gọi bị chặn với \texttt{duration=0ms}, nghĩa là backend không hề nhận được yêu cầu. Thứ ba, mỗi lệnh gọi để lại dấu vết ở nhiều phía độc lập với nhau, gồm phép kiểm ở client, thông tin mock báo lại ở backend, access log ở gateway và dashboard. Các phía này khớp nhau. Và chính vì độc lập, một phía ghi sai sẽ tự lộ ra qua chênh lệch với những phía còn lại.

\subsubsection{Hạn chế}

Hạn chế lớn nhất là bằng chứng và cấu hình hiện tại không hoàn toàn trùng nhau. Toàn bộ bằng chứng được ghi khi năm target đều là mock và các chính sách chỉ chứa quy tắc của agent. Từ đó đến nay, cấu hình đã gộp quy tắc người dùng và agent vào một chính sách, chuyển radar, Spark History và OpenMetadata sang phần mềm thật, lấy \texttt{x-forwarded-user} từ \texttt{preferred\_username} và áp NetworkPolicy cho mọi pod trong namespace backend, nhưng chưa có đầu ra kiểm thử đầu cuối nào trên cấu hình hợp nhất này. Mock lại không gây tác dụng phụ, nên kết quả chưa nói được gì về quyền trên tài nguyên thật như namespace, DAG hay thực thể metadata. Việc đổi token cho OpenMetadata cũng chưa có bằng chứng.

Độ phủ ở bước gọi còn mỏng. 210 ô được so ở bước khám phá, song chỉ 18 lệnh được gọi thật, và chỉ một lệnh thử tách đọc khỏi ghi trên cùng hệ thống. Bộ kiểm thử coi mọi ngoại lệ phía client là ``đã bị chặn'', nên tự nó không phân biệt được một lệnh bị từ chối đúng cách với một lỗi kết nối; mã 400 được xác nhận qua access log và dashboard chứ không qua bộ kiểm thử. Dashboard cũng đếm gộp lệnh thiếu quyền với lệnh gọi sai tên công cụ, hệ quả của việc gateway trả cùng một thông báo \texttt{Unknown tool} cho cả hai. Về vận hành, số đo khởi động mới có năm lần ở tình huống khởi động ấm. Thời gian dựng từ cụm trống và độ trễ mà gateway thêm vào mỗi lệnh chưa được đo; lớp bọc PowerShell chưa được chạy trực tiếp vì phiên đo chạy bên trong WSL; script dọn dẹp còn để nguyên hai nhóm container compose; và tracing chưa được cấu hình.

\endgroup
